\section{Implementation, independent validation, and measured performance}
\label{main:experiments}

The delivered implementation adds a certified terminal operation on a supplied
normal cut, together with the separate sector decomposition of
Section~\ref{sec:sector-blocks}.  Its computational evidence addresses three
different questions: whether the certificates survive independent replay,
which supplied cuts satisfy the current sufficient criterion, and how the
cost of the compressed operation compares with constructing and recognizing
an explicit cut.  The experiments do not measure a complete run from a knot
diagram to a knot verdict.

\subsection{Additive modules and verification boundaries}
\label{main:exp-implementation}

The contribution consists of eight new modules under \texttt{fast/fastunknot/},
arranged as four producer/checker pairs.  Existing recognition defaults and
existing production modules are unchanged.

\begin{table}[htbp]
\centering
\small
\begin{tabularx}{\linewidth}{@{}p{0.23\linewidth}Y@{}}
\toprule
Operation & New modules \\
\midrule
Ordered interfaces &
\texttt{parallelity\_boundary\_words.py}\par
\texttt{parallelity\_boundary\_words\_verify.py} \\
Geometric core cancellation &
\texttt{core\_handlebody.py}\par
\texttt{core\_handlebody\_verify.py} \\
Whole-cut census &
\texttt{normal\_cut\_handlebodies.py}\par
\texttt{normal\_cut\_handlebodies\_verify.py} \\
Sector factors &
\texttt{sector\_blocks.py}\par
\texttt{sector\_blocks\_verify.py} \\
\bottomrule
\end{tabularx}
\caption{Additive implementations and their separate certificate checkers.
The research drivers, fixtures, and tests are additional artifacts.}
\label{main:exp-modules}
\end{table}

The interface checker reconstructs its geometry through the existing
raw-cut-derived reference implementation.  It builds the two orbit systems
independently, invokes the established weighted-orbit checker, and reconstructs
the signed words and maximal attaching intervals without calling the new
producer's query builder or traversal.  The core checker obtains unsigned
incidences from authenticated handle cycles; it accepts no supplied incidence
summary.  It verifies the spanning forest and either the immutable leaf order
or the saturated geometric pivot sequence.  The whole-cut checker first
authenticates both source decompositions and then independently recomputes the
attachment graph, its connected components, and its genus histogram.

This independence is an implementation boundary, not a claim that the checkers
have no shared foundations.  They deliberately reuse established source
validation, exact integer transport, and orbit-proof verification.  The
geometric arguments in Sections~\ref{int:section} and~\ref{geo:section} justify
the interpretation of the checked data.  Tests that disable the producer
functions ensure that certificate acceptance does not accidentally call the
discovery routines.

The interface schema is \texttt{parallelity-boundary-words-v1}; the conditional
core schema is \texttt{normal-core-handlebody-v1}; and the source-bound whole-cut
schema is \texttt{normal-cut-handlebodies-v1}.  For the last of these,
\texttt{COMPLETE} asserts the exact binary multiplicity census of the supplied
cut.  \texttt{INCONCLUSIVE} asserts neither non-handlebody topology nor
knottedness.  A core proof by itself remains conditional on the authenticity
of its geometric handle charts.

The production allowance \texttt{max\_cycles} is shared by the four inherited
parallelity queries and the two new interface queries.  Exhausting it leaves
the attempt unresolved, without a partial accepted census.  The checker can
cap the total number of supplied orbit operations through
\texttt{max\_operations}; cooperative cancellation also reaches the finite
core and graph computations.  These are work controls, not wall-clock
guarantees.  Large integers use the repository's exact hexadecimal JSON
transport where needed.  The separate sector schemas and their more limited
Euler and ray-enumeration meanings are described in
Section~\ref{sec:sector-blocks}; none silently asserts an essential disc.

\subsection{Regression tests and adversarial checks}
\label{main:exp-tests}

The recorded integrated run passed \textbf{1,644 tests in 346.562 seconds};
the log is \texttt{integrated-tests.txt}.  Focused validation also passed the
13 core/census methods in \texttt{core-census-tests.txt} and the eight boundary
methods in \texttt{boundary-words-tests.txt}.  These focused counts overlap
the integrated suite and are not added to it.

The tests check source relabellings, one-sided normal surfaces, empty cuts,
parallel meridians, interior-vertex obstructions, repeated incidences, exact
large integers, shared budgets, cancellation, and producer-disabled replay.
Adversarial mutations alter source fingerprints, endpoint tags, fibre
transport, signed words, attachment types, core owners, forest labels,
cancellation schedules, and final genera.  The expected outcome is rejection
of the altered certificate.

Two finite algebraic audits target the central cancellation mechanisms.
For 1,200 deterministic random incidence systems, every produced saturated
schedule is replayed against unbounded nonnegative integer counts.  For
750 deterministic random square matrices of orders one through five, the
leaf producer's success is compared with a directly computed permanent.
The latter test exposed a stale-queue bug during development: a row already
queued as a leaf could cease to be eligible after another column was removed.
The producer now rechecks both row survival and its current unique unit
entry when popping the queue.  This correction preceded the final passing
runs.  The permanent computation is a test oracle for these small matrices,
not an operation used by the classifier.

The boundary tests cover 193 source/vector cases, of which 182 are compared
with a literal expanded polygon-boundary oracle.  That oracle walks actual
boundary edges and reconstructs the corner/fibre identifications separately
from the compressed word routine.  Eleven large cases use independent
compressed replay and source consistency checks instead of exceeding that
oracle's expansion guard.  Section~\ref{int:evidence} gives the breakdown and
the separate interface-scale measurements.  The nine focused sector test
methods, their dense independent ray comparisons, the 188-sector audit, and
the sector timing experiment are reported in Section~\ref{sec:sector-blocks}.

\subsection{A source-cut census with an independent topological oracle}
\label{main:exp-census}

The file \path{census-audit.json} records all 106 supplied cuts in
\path{marked_parallelity_research/corpus.json}, tied to source revision
\texttt{d2c5e1707fe783bbed06743ac732d70d83e6f27b}.  Every native success was
independently replayed.  For manageable coordinates, a separate Regina path
constructs the triangulation from the source face pairings, constructs the
normal surface, calls \texttt{cutAlong()}, separates the resulting components,
and calls \texttt{recogniseHandlebody()} on each component.  It uses no native
core, attachment graph, or genus calculation.

\begin{table}[htbp]
\centering
\small
\begin{tabular}{@{}lr@{}}
\toprule
Recorded outcome & Cases \\
\midrule
Native \texttt{COMPLETE}, with successful independent replay & 30 \\
Native \texttt{INCONCLUSIVE} & 76 \\
\quad Full-circle core attachment & 60 \\
\quad Raw cancellation stalled & 6 \\
\quad Core has three-handles & 10 \\
\midrule
Explicit Regina oracle completed & 99 \\
Native successes compared with Regina; all agreed & 23 \\
Binary cases above the explicit expansion cap & 7 \\
\bottomrule
\end{tabular}
\caption{Whole-cut audit.  Oracle rows describe a second view of the same
106 cases and are not additional cases.  All seven expansion-capped cases
have successful native certificates.}
\label{main:exp-census-counts}
\end{table}

The corpus oracle refuses expansion above 4,096 normal disks, before
constructing a Regina cut; its subprocess also has a per-case timeout.
The recorded run completed all 99 uncapped oracle computations and encountered
no oracle timeouts.  The seven binary cases were explicitly marked
\texttt{NOT\_ATTEMPTED\_EXPANSION\_CAP}.  They were checked through native
certificate replay, without attempting an overflowing or unrepresentable
explicit cut and without counting the skip as an oracle agreement.

The 30 successful censuses and 23 independent topological agreements are
finite implementation evidence for the proved criterion.  The 76 unresolved
cases show that the criterion remains restrictive on this selected corpus.
In particular, a circle obstruction is a failure of the delivered
disk-interface consumer, not a claim that the actual cut is not a union of
handlebodies.

The separate \texttt{annular-absorption-diagnostic.json} examines precisely
those 60 circle-obstructed cases in light of
Proposition~\ref{geo:prop-annulus-absorption}.  In 38 cases every obstructing
full-circle attachment belongs to a disk base; the other 22 have at least
one circle on a non-disc base.  Fifteen additional boundary queries needed
by the diagnostic were independently replayed.  This is an eligibility
diagnostic: it did not append the proposed new two-handles, run augmented
cancellation, or accept an additional census.  Consequently the measured
coverage remains 30 of 106, rather than 68 of 106.

\subsection{Paired timings of the complete supplied-cut operation}
\label{main:exp-timing}

The benchmark environment recorded in \texttt{census-benchmark.json} is
Python 3.12.14, Regina 7.4, and Linux on x86-64.  Timings use
\texttt{time.perf\_counter}.  Both arms start from the same face pairings and
the same supplied standard normal vector in a Fibonacci layered solid torus.
The vector is the primitive meridian vector multiplied by $m$.  Every timed
case is checked to give exactly $m$ ball components.

The native arm includes certificate production and independent source-bound
replay.  The Regina arm includes construction of its triangulation and normal
surface, explicit cutting, component separation, and handlebody recognition.
JSON serialization and hashing are outside the native timer.  Each case has
an initial correctness/warm-up execution of both arms, followed by five
paired timing rounds.  The order within each pair is shuffled using the
recorded seed 20261012.

For round $r$, write $n_r$ for native production plus replay and $e_r$ for the
explicit Regina operation.  The two time columns below are their separate
medians.  The comparison statistic is
\[
 R=\operatorname{median}_{r=1,\ldots,5}\frac{e_r}{n_r},
 \qquad
 R\ne
 \frac{\operatorname{median}_r e_r}{\operatorname{median}_r n_r}
 \quad\hbox{in general}.
\]
Thus $R>1$ favors the native arm.  Each case also has three additional
Regina-versus-Regina pairs; the final column gives the median of their
first-run/second-run ratios.  These A/A controls describe timing variability,
not a second algorithmic comparison.

\begin{table}[htbp]
\centering
\small
\setlength{\tabcolsep}{3.4pt}
\begin{tabular}{@{}rrrrrrrr@{}}
\toprule
$T$ & $m$ & $D$ & $T_{\rm cut}$ &
\shortstack{Native\\median (s)} &
\shortstack{Regina\\median (s)} &
$R$ & A/A \\
\midrule
1 & 1 & 3 & 31 & 0.001794 & 0.000341 & 0.205 & 1.128 \\
1 & 16 & 48 & 226 & 0.006797 & 0.001470 & 0.217 & 1.051 \\
1 & 256 & 768 & 3,346 & 0.007451 & 0.022121 & 2.922 & 0.965 \\
1 & 4,096 & 12,288 & 53,266 & 0.007727 & 0.429135 & 54.974 & 0.989 \\
4 & 1 & 29 & 208 & 0.183885 & 0.005508 & 0.030 & 1.064 \\
4 & 8 & 232 & 1,013 & 0.486196 & 0.026918 & 0.054 & 0.980 \\
4 & 64 & 1,856 & 7,453 & 0.458985 & 0.182686 & 0.394 & 0.968 \\
8 & 1 & 228 & 1,093 & 1.665820 & 0.171979 & 0.105 & 1.009 \\
8 & 8 & 1,824 & 7,393 & 2.325465 & 1.145137 & 0.499 & 0.994 \\
12 & 1 & 1,592 & 6,573 & 4.704264 & 8.618704 & 1.817 & 1.058 \\
\bottomrule
\end{tabular}
\caption{All ten paired whole-cut benchmark cases.  $D$ counts normal disks;
$T_{\rm cut}$ is the number of tetrahedra in Regina's explicit cut.
Each arm median and the median paired ratio are computed from the saved
individual samples.}
\label{main:exp-timing-table}
\end{table}

The results show both overhead and useful savings.  Regina is faster in seven
of these ten cases, including the small one-copy inputs at $T=1,4,8$.
The native operation becomes favorable on the one-tetrahedron repeated-disc
examples at $m=256$ and $m=4096$, and on the $T=12$, $m=1$ input.  The largest
recorded paired ratio is about $54.97$ at $(T,m)=(1,4096)$; its explicit cut
contains 53,266 tetrahedra.  At $(12,1)$, the corresponding ratio is about
$1.82$, with both native proof phases included.

These finite comparisons support the practical value of avoiding explicit
normal-stack expansion on suitable inputs.  They do not assert a universal
speed advantage.  Five samples per arm are descriptive measurements; no CPU
affinity, isolated memory environment, or peak-memory measurement was
established for this run.  The saved A/A data also include noticeable
short-duration variation.  No whole-recognizer runtime or empirical
quasi-polynomial exponent is inferred from this table.

The benchmark has a separate explicit limit of 50,000 normal disks.  It is
deliberately distinct from the audit's 4,096-disk guard: the largest timed
explicit example has $D=12,288$ and is allowed here.  Neither limit is a
bound on the number of tetrahedra in the constructed cut.

\subsection{Binary capacity and actual certificate sizes}
\label{main:exp-capacity}

Four additional measurements run only the native producer and independent
checker, with multiplicity $m=2^p$.  All return and verify the histogram
$\{(0,2^p)\}$.  They are capacity measurements, each from one execution;
there is no explicit-oracle timing or speed ratio for these enormous cuts.
The serialized certificates are measured as sorted, uncompressed JSON using
the project's exact integer codec.  Decimal megabytes mean $10^6$ bytes.

\begin{table}[htbp]
\centering
\small
\setlength{\tabcolsep}{3.2pt}
\begin{tabular}{@{}rrrrrrrr@{}}
\toprule
$T$ & $p$ &
\shortstack{Produce\\(s)} &
\shortstack{Replay\\(s)} &
\shortstack{Total\\(s)} &
\shortstack{Certificate\\bytes} &
\shortstack{Decimal\\MB} & $P/a$ \\
\midrule
1 & 16,384 & 0.008944 & 0.011613 & 0.020556 &
1,701,021 & 1.701021 & 0/0 \\
4 & 1,024 & 0.360605 & 0.241919 & 0.602524 &
2,510,984 & 2.510984 & 15/3 \\
4 & 16,384 & 1.224719 & 0.858719 & 2.083438 &
30,078,776 & 30.078776 & 15/3 \\
8 & 16,384 & 6.404215 & 4.195131 & 10.599347 &
71,291,795 & 71.291795 & 47/3 \\
\bottomrule
\end{tabular}
\caption{Separate native binary-capacity measurements.  The last column
lists ports $P$ and maximal interval attachments $a$, rather than their
numerical quotient.  Serialization and hashing are outside the reported
production/replay times.}
\label{main:exp-capacity-table}
\end{table}

The fixed number of exceptional interfaces does not imply a constant-sized
proof.  For example, the eight-tetrahedron case still has 47 ports and three
interval attachments at $m=2^{16384}$, but its full census certificate occupies
71.291795 decimal MB.  Source coordinates and binary interval endpoints recur
through the recorded proof traces.  These are genuine output costs even
though the $2^{16384}$ ball components are never individually materialized.
The recorded SHA-256 digests identify the exact serialized proofs used in
the capacity measurements.

\subsection{Reproduction and the limits of the evidence}
\label{main:exp-reproduction}

The complete per-case census is in \texttt{census-audit.json}; paired samples,
arm order, A/A samples, capacity times, certificate lengths, and certificate
digests are in \texttt{census-benchmark.json}.  The diagnostic retains its
per-case eligibility decisions in
\path{annular-absorption-diagnostic.json}.  From the repository's
\texttt{fast/} directory, the supplied census drivers are invoked as follows:
\begin{lstlisting}
python -m marked_parallelity_research.audit_census \
    --output census-audit.json
python -m marked_parallelity_research.benchmark_census \
    --output census-benchmark.json --rounds 5 --capacity
\end{lstlisting}
Regina is required for the explicit comparison arm.  Native certificate
production and replay use the repository's Python implementation.

The experimental conclusions have deliberately separate scopes.  Passing
tests and successful independent replay give evidence about implementation
correctness.  Agreement with independently constructed Regina cuts checks
the topological interpretation on the manageable part of the corpus.
Capacity runs demonstrate processing of large binary multiplicities and
measure the resulting proof sizes.  The all-input geometric and asymptotic
claims rest on the theorems, with their stated source and interface
hypotheses.  Surface discovery, peripheral-data transport, complete hierarchy
coverage, and a bound on the global recognition search remain outside these
timings and outside the measured success count.
